\subsection{3. Mục tiêu đề tài}
\indent\indent Mục tiêu của đề tài là nghiên cứu và xây dựng hệ thống phân loại video, hình ảnh, văn bản và âm thanh văn hoá phi vật thể bằng việc tận dụng mô hình đồ thị đa phương thức kết hợp với mạng KAN. Đề tài hướng tới việc tối ưu hóa hiệu suất phân lớp bằng cách cho mô hình học các mối quan hệ phức tạp giữa thông tin hình ảnh, văn bản (phụ đề) và âm thanh. Kết quả nghiên cứu sẽ cung cấp cơ sở thực nghiệm cho việc ứng dụng các kiến trúc mạng noron thế hệ mới (KAN) vào bài toán phân lớp dùng đồ thị đa phương thức. Các kết quả cụ thể mà đề tài hướng đến bao gồm:

\begin{itemize}[label={--}, itemsep=1pt]
    \item \textbf{Thu thập và tiền xử lý dữ liệu đa phương thức:} Thu thập dữ liệu video và phụ đề tương ứng từ nền tảng Youtube thuộc các chủ đề văn hóa đặc trưng như: Vía bà chúa xứ Núi Sam, Nghề đan tre, Nghề dệt chiếu, Lễ hội nghinh Ông,... Thực hiện các kỹ thuật tiền xử lý bao gồm: cắt video thành các đoạn ngắn (video segments), đồng bộ hóa phụ đề, và lọc bỏ những mẫu dữ liệu nhiễu hoặc không liên quan đến nội dung văn hóa để xây dựng bộ 
    dữ liệu huấn luyện và kiểm thử chuẩn.
    \item \textbf{Xây dựng và đánh giá mô hình đơn phương thức (baseline):} Xây dựng mô hình baseline chỉ để phân loại chỉ sử dụng dữ liệu hình ảnh. Sử dụng các kiến trúc mạng noron tích chập (CNN) phổ biến đã được huấn luyện trên tập ImageNet (như ResNet, RegNet,...) làm bộ trích xuất đặc trưng nền tảng (backbone). Tiến hành đóng băng (freeze) các trọng số của backbone và huấn luyện tầng phân loại (classifier) để đánh giá hiệu quả khi chỉ dùng hình ảnh. Kết quả này đóng vai trò là mốc so sánh và là thành phần trích xuất đặc trưng cho mô hình đồ thị đa phương thức.
    \item \textbf{Xây dựng và đánh giá mô hình MGNN kết hợp KAN:} Xây dựng mô hình đồ thị đa phương thức sử dụng mạng noron đồ thị (GNN), cụ thể hơn là Graph Attention Network (GAT). Tận dụng backbone hình ảnh đã thử nghiệm ở giai đoạn (2) để tận dụng đặc trưng hình ảnh giàu thông tin. Xây dựng các tầng GAT để học mối quan hệ ngữ nghĩa giữa hình ảnh, văn bản, và âm thanh. Đồng thời tích hợp bộ phân lớp bằng KAN để tăng tính biểu diễn phi tuyến, với kỳ vọng đạt được độ chính xác cao hơn, ngoài ra tôi cũng thử trên mạng MLP và mô hình máy học truyền thống để có cái nhìn so sánh trực quan.
    \item \textbf{Triển khai hệ thống trên nền tảng web:} Xây dựng ứng dụng web trực quan để minh họa kết quả nghiên cứu. Hệ thống cho phép người dùng tải lên video, hình ảnh, văn bản và âm thanh. Đầu ra của hệ thống là nhãn chủ đề văn hóa được dự đoán bởi mô hình, kèm theo độ tin cậy của dự đoán, giúp người dùng dễ dàng kiểm chứng hiệu quả của mô hình.
\end{itemize}